% =====================================================================
%  第 4 章  阶段结果
% =====================================================================

\section{阶段结果}
\label{sec:results}

本章报告四组实验。前三组不涉及大语言模型，
用于确认数据、线性规划与统计筛选三部分各自正确；
第四组给出智能体接入模型后的完整运行结果，并以消融实验检验模型的实际贡献。

\subsection{数据的整体面貌}

数据集共 27810 张图。图\ref{fig:dist} 给出五个不变量在各阶内部的分布比例。
随阶数增大，色数与团数的分布逐渐分开，独立数向中段集中。

\begin{figure}[H]
\centering
\includegraphics[width=0.92\linewidth]{fig01_invariant_dist.png}
\caption{各不变量在每一阶内部的分布比例}
\label{fig:dist}
\end{figure}

图\ref{fig:corr} 给出九个不变量两两之间的秩相关系数。
色数与团数的相关系数为 0.97，边数与谱半径为 0.96，
独立数与边数为 $-0.43$。前两者说明这些量高度冗余，
后者说明独立数与图的稠密程度反向相关。
这一观察直接决定了候选形式的选择范围：
若两个量几乎完全相关，则同时使用二者不提供额外信息。

\begin{figure}[H]
\centering
\includegraphics[width=0.78\linewidth]{fig03_corr_heatmap.png}
\caption{不变量两两之间的秩相关系数}
\label{fig:corr}
\end{figure}

\subsection{线性规划能否复现已知结论}

线性规划的正确性由一项检验确认：给定已知结论所对应的形式，
规划求出的系数应当复现该结论。

图\ref{fig:recovery} 给出五条已知结论在 $n \leq 8$ 的全部 13598 张图上的取等比例。
五条均被正确复现，其中 $\chi \leq \omega+1$ 与 Reed 型不等式各在 4.3\% 的图上取等，
$\chi \leq n - \alpha + 1$ 与 $\alpha \leq n - \omega + 1$ 各在 3.1\% 的图上取等，
$\chi \leq \Delta + 1$ 的取等比例最低，为 0.7\%。

\begin{figure}[H]
\centering
\includegraphics[width=0.92\linewidth]{fig06_lp_recovery.png}
\caption{五条已知结论在数据上的取等比例}
\label{fig:recovery}
\end{figure}

取等比例随阶数下降。在三阶图上，三条结论的取等比例均为 100\%；
到八阶时降至 1\% 至 2\%。这说明取等只在少数结构特殊的图上发生，
大多数图上的不等式是严格成立的。
这一点对后续的统计筛选有直接影响：取等比例本身就很低，
故不能以“取等比例高”作为候选有效的唯一依据，
还需与零分布比较。

\subsection{小阶成立不蕴含普遍成立}

图\ref{fig:boundary} 给出一个值得单独说明的现象。
不等式 $\chi \leq \omega + 1$ 在 $n \leq 8$ 的全部 13598 张图上成立，
在 $n = 9$ 与 $n = 10$ 的采样图上亦成立，
却在 $n = 11$ 上失效。反例是 Grötzsch 图：
11 个顶点、20 条边、团数为 2、色数为 4。
该图不含三角形，故团数仅为 2，但需要四种颜色。

\begin{figure}[H]
\centering
\includegraphics[width=0.95\linewidth]{fig09_boundary.png}
\caption{$\chi \leq \omega+1$ 在每一阶上的最大超出量}
\label{fig:boundary}
\end{figure}

这一现象是本课题核心判断的直接证据。
若仅在小规模图上验证，$\chi \leq \omega+1$ 看起来是一条可靠的规律，
甚至可能被当作新发现报告；但它并不普遍成立。
因此系统必须把“在数据上成立”与“成立”严格区分，
而统计筛选正是为此设置的。

值得注意的是，该反例由结构化族在预算 45 内找到，
而随机图搜索在同一预算下几乎无法命中。
这为“搜索策略需要判断”提供了实测依据。

\subsection{反例搜索策略的比较}

对五条候选，各用三类策略作八次独立搜索，每次预算 2000 张图。
结果如表\ref{tab:strategy} 所示。

\begin{table}[H]
\centering
\small
\caption{三类搜索策略的命中情况（成功次数 / 8，括号内为命中时的中位预算）}
\label{tab:strategy}
\begin{tabular}{@{}lccc@{}}
\toprule
候选 & 随机图 & 结构化族 & 局部修改 \\
\midrule
$\chi \leq \Delta$ & 8（5） & 8（1） & 8（3） \\
$\chi \leq \omega + 1$ & 0 & 8（50） & 0 \\
$\alpha + \omega \leq n$ & 8（3） & 8（5） & 8（1） \\
$\chi \leq \rho + 1$ & 0 & 0 & 0 \\
$\chi \leq \frac{1}{2}(\omega+\Delta) + 1$ & 0 & 0 & 0 \\
\bottomrule
\end{tabular}
\end{table}

三点结论。其一，两类策略各有其适用面。
结构化族是唯一能推翻 $\chi \leq \omega+1$ 的策略，
八次全部命中，中位预算 50；随机图八次全部落空。
原因在于该命题的反例必须是无三角形而色数不小于 4 的图，
随机撒点几乎不可能撞上这种结构，
而结构化族里的无三角形构造正是为此准备的。

其二，局部修改的表现恰好相反。
它在 $\chi \leq \Delta$ 与 $\alpha+\omega \leq n$ 上最快，
中位预算分别是 3 与 1，远低于另外两类；
但对 $\chi \leq \omega+1$ 八次全部落空。
局部修改从“取等图”出发做小步扰动，
而 $\chi \leq \omega+1$ 的反例与取等图相距过远，
扰动无法跨越这段距离。这正说明三类策略不是互相替代的关系：
它们的有效区间由候选反例的几何位置决定，
而非由预算大小决定。

其三，后两条候选在三种策略下均未找到反例。
$\chi \leq \rho+1$ 是 Wilf 定理的特例，
$\chi \leq \frac{1}{2}(\omega+\Delta)+1$ 则是 Reed 猜想的松弛形式。
后者自 1998 年提出至今未解决，
若本系统能在数秒内将其推翻，反倒说明实现有误。
这一结果与理论预期一致，可作为系统未产生荒谬结论的旁证。

表\ref{tab:strategy} 最后一行需要一处说明。
若让线性规划自行拟合 $\chi \leq \frac{1}{2}(\omega+\Delta)+1$ 的形式，
它会将 $\Delta$ 的系数压成零，退化为 $\chi \leq \omega+1$，
于是“命中”实为另一条命题的反例。
表中的结果因此是手工指定系数 $0.5$ 与 $0.5$ 后得到的。
这一细节本身也是前文所说“塌缩”的一种表现：
形式空间中的两个元素在拟合后落到同一条不等式上，
若不检查系数是否为零，就会把两条命题混为一谈。

图\ref{fig:refute} 以中位预算为横轴展示了前三条候选的结果，
未命中的组合在图中不出现。
表中后两条候选在三种策略下均未命中，故不在图中。

\begin{figure}[H]
\centering
\includegraphics[width=0.95\linewidth]{fig07_refute_curve.png}
\caption{三类策略命中反例所需的中位预算}
\label{fig:refute}
\end{figure}

\subsection{反例搜索的可靠性}
\label{sec:injection}

上一小节比较的是策略之间的相对优劣。
但真实候选的真伪无从判定，因此无法直接回答一个更要紧的问题：
\textbf{当一条命题只在极少数图上失效时，搜索还找得到吗？}
找不到反例时，系统会把候选送进统计筛选，
若这一步漏掉了本该被推翻的命题，后面的筛选就建立在错误的候选上。

为此设计了一组注入实验：先用已知为假的命题替换真实候选，
再测量各类策略能否找到反例。
难点在于构造合适的假命题。最初的做法是把已知成立的界整体收紧一或二，
但九个不变量都取整数值，松弛量大量并列，
收紧一格就从零个反例跳到成百上千个，
三类策略在最小预算下检出率即达百分之百。
这测的是“能否发现一条明显错误的命题”，与真实困难无关。

最终改为从数据中筛出天然的难例：
在完整形式空间上逐条求解，保留在整个数据集上
\textbf{违反数不超过 5 张、且这些图的阶数都不小于 9} 的形式，共得 30 条。
这类命题在绝大多数图上成立，只在少数大阶图上失效，
正是真实场景中的困难情形。
对每条命题、每种策略、每档预算各重复三次，结果如图\ref{fig:injection} 所示。

\begin{figure}[H]
\centering
\includegraphics[width=0.98\linewidth]{fig14_injection.png}
\caption{注入难例上的检出率（左：随预算变化；右：按反例张数分档）}
\label{fig:injection}
\end{figure}

结果并不乐观。预算从 50 加到 3000，即扩大六十倍，
结构化族的检出率只从 0.27 升到 0.30，局部修改恒为 0.10，随机图恒为零。
按反例张数分档更说明问题：
整份数据集上只有 1 张反例的 18 条命题，检出率仅 \textbf{0.02}；
有 2 张与 3 张反例的，检出率也只有 0.34 与 0.33。

两点解释。其一，检出率几乎不随预算变化，
说明瓶颈不在算力而在\textbf{搜索分布的支撑集}。
结构化族之所以能把检出率从零提到 0.27，
是因为它的候选池里恰好含有无三角形、高色数这类构造；
一旦池中的确定性部分用尽（实测 122 张），
后续补入的只是随机图，检出率便不再上升。
其二，反例张数越少越难找，
与“反例是测度零的集合”这一直觉一致：
随机撒点命中概率正比于反例所占比例，
而该比例在难例上低至万分之一量级。

这组实验给出了本系统最重要的一条限制：
\textbf{“未找到反例”在难例上几乎不提供信息}。
第\ref{sec:ceiling}节的上界分析会显示，
在完整形式空间里能通过全部检验的候选有三十余条，
而智能体五轮只提出四十条形式，
两者能对上的原因并非搜索充分，
而是绝大多数形式本身就不成立、在第一关就被淘汰。
对于真正接近成立的那些候选，
现有预算下的反例搜索无法把它们与成立的命题区分开。

\subsection{统计筛选的效果}

在不调用模型的基线上，完整跑通生成与筛选两步。
目标不变量取色数，形式空间限定为至多两个预测子，
共 128 种形式。流程与各级淘汰数如图\ref{fig:funnel} 所示。

\begin{figure}[H]
\centering
\includegraphics[width=0.88\linewidth]{fig10_funnel.png}
\caption{候选的逐级淘汰}
\label{fig:funnel}
\end{figure}

128 种形式中，52 种求得了非平凡解，其余因形式过弱而退化为平凡不等式。
去重后剩 17 条。经 200 次置换检验并作 BH 校正（$q = 0.1$）后保留 7 条。
最后在从未参与拟合与筛选的测试集上复核，7 条中有 6 条仍然成立。

保留的六条中有四条是已知结论的复现：
$\chi \leq \omega+1$、$\chi \leq n-\alpha+1$、$\chi \leq \rho+1$、$\chi \leq \Delta+1$。
另外两条为带小数的组合形式，例如
$\chi \leq 0.4\,\delta + 0.6\,\nu + 4.2$。
这类形式的取等比例极低（0.06\%），
其数学意义尚待判断，不宜径直称为新发现。

唯一未通过测试集复核的是 $\chi \leq 0.667\,\Delta + 0.125\,\delta + 2.67$。
该式在发现集与验证集上均成立，在测试集上失效，
说明它是一条过拟合的产物。这一条恰好说明测试集的设置是必要的：
若只作发现集与验证集的两分，该式会被当作有效候选保留下来。

图\ref{fig:pvalue} 给出 17 条候选的 $p$ 值分布与 BH 过程。
被保留的 7 条 $p$ 值均为 0.005，即 200 次置换中零分布从未达到观测值，
这是 $B=200$ 时 $p$ 值可能取到的最小值。

\begin{figure}[H]
\centering
\includegraphics[width=0.95\linewidth]{fig08_pvalue_hist.png}
\caption{候选的 $p$ 值分布与 BH 校正过程}
\label{fig:pvalue}
\end{figure}

\subsection{智能体接入模型后的运行}
\label{sec:agentrun}

将大语言模型接入主循环后，完整运行五轮。最终一次运行共 26 次模型调用，
消耗输入 16864 个、输出 5411 个词元。每轮提出 8 条形式，五轮共 40 条。
各轮的去留情况如图\ref{fig:agentfunnel} 所示。

\begin{figure}[H]
\centering
\includegraphics[width=0.96\linewidth]{fig11_agent_funnel.png}
\caption{一次完整运行中候选的逐轮去留与累计}
\label{fig:agentfunnel}
\end{figure}

五轮共推翻 5 条，保留 3 条，三条在测试集上全部成立：

\begin{enumerate}[leftmargin=2.2em,itemsep=0.15em]
\item $\alpha \leq n - \nu$，取等比例 0.3158；
\item $\alpha \leq n - \omega + 1$，取等比例 0.0177；
\item $\chi \leq \rho + 1$，取等比例 0.0046。
\end{enumerate}

三条候选的 $p$ 值均为 0.0050，即 200 次置换中零分布从未达到观测值。
事后核对，三条其实都是已知结论。$\alpha \leq n - \nu$ 即
“最大独立集与最大匹配的顶点不相交”，$\alpha \leq n - \omega + 1$ 是初等界，
$\chi \leq \rho + 1$ 则是 Wilf 在 1967 年证明的定理：
色数不超过 $1 + \lambda_{\max}$。邻接矩阵非负，由 Perron--Frobenius 定理，
其谱半径本身就是一个特征值，故 $\rho = \lambda_{\max}$，两条式子完全一致。
系统独立地重新发现了这条定理，这是“提出—反驳—筛选”链路能够复现已知结果的一个例证。

这一核对结果值得单独说明。系统在三条候选上都没有产出新结论，
但三条都对应真实的定理，说明“提出—反驳—筛选”这条链路
确实能复现已知结果。同时它也暴露了新颖性判定的脆弱：
第三条在初次核对时被判为“查不到出处”，
补查特征值与色数的经典结果之后才归入已知。
检索不到与不存在是两回事，而系统只能做前者。
模型给出的理由同样不可靠：它对第一条称“Gallai 恒等式，二分图情形”，
而该不等式在一般图上就成立。这说明新颖性判定即使结论正确，
依据也可能站不住，只能作为人工复核的线索，不能单独作为筛选依据。

值得一并说明的是，早前一次运行还产出过 $\rho \geq \omega - 1$。
该条在测试集上成立，初看也像新结论，实则是初等的：
取最大团的指示向量作 Rayleigh 商，即得 $\rho \geq \omega - 1$。
它在后续运行中没有再被提出，
说明智能体的产出对提示词的细微变化敏感，
同一次运行的保留集合不能视为该系统的稳定能力边界。

\subsection{三轮迭代：反馈信息如何影响产出}

上述结果并非一次运行得到。系统在调试中经过三次修改，
每一轮修改都针对轨迹中暴露的一个具体缺陷，改动与效果见表\ref{tab:iterations}。

\begin{table}[H]
\centering
\small
\caption{三轮迭代的改动与效果}
\label{tab:iterations}
\begin{tabularx}{\linewidth}{@{}c
  >{\raggedright\arraybackslash}p{5.0cm}
  >{\raggedright\arraybackslash}p{2.6cm}
  >{\raggedright\arraybackslash}X@{}}
\toprule
版本 & 针对的缺陷与改动 & 调用次数 & 结果 \\
\midrule
v4 & 基线：每轮只告知“哪些形式试过了” & 47 & 保留 4 条 \\
v5 & 把反例的数值与模型自己归纳的失败原因回传下一轮
   & 35 & 保留 2 条，更差 \\
v6 & 追加两项反馈：被线性规划压成零的系数、已提形式的键
   & 35 & 保留 4 条，重复大幅下降 \\
v7 & 修正恒真式识别：补上同一不等式的对称写法（第\ref{sec:ceiling}节）
   & 26 & 保留 3 条，调用最少 \\
\bottomrule
\end{tabularx}
\end{table}

\textbf{v5 的失败值得单独说明。}直觉上，把失败原因告诉模型应当有益，
但实测保留条数从 4 降到 2。检查轨迹后确认原因：
提示词中列出的“已试形式”用的是不等式文本（如 $\chi \leq \rho + 1$），
而模型被要求输出的是形式（如 \texttt{chi le [+rho]}），
两者字面不同，模型无法把二者对应起来，于是反复提出同一形式。
加入失败原因只是让提示词变长，并未解决这一对应问题。
v5 的第 3 至 5 轮分别有 5、7、6 条形式与前轮重复。

\textbf{v6 的改动针对的正是这个对应问题。}轨迹中另有两类信息此前被丢弃：
一是模型提出三个预测子、线性规划把其中一两个系数压到零后塌缩为旧形式，
模型看不到自己的形式为何“消失”；二是恒真式被剔除后模型不知道哪些组合无意义。
v6 把这两类信息与已提形式的键一并写回提示词。
结果重复显著下降，第 4 轮首次出现 8 条提议全部塌缩为已检验过的不等式的情形。

图\ref{fig:trace} 给出 v7 每轮 8 条提议的去向。

\begin{figure}[H]
\centering
\includegraphics[width=0.9\linewidth]{fig13_trace_timeline.png}
\caption{各轮提议的去向构成}
\label{fig:trace}
\end{figure}

第 2 至 5 轮分别有 3、2、3、3 条形式与前轮重复，
第 4 轮 8 条提议全部塌缩或重复，没有一条产生新候选；
第 5 轮只有 1 条可解，但正是这一条被保留。
这说明即使反馈信息完整，模型在五轮之后也接近耗尽它在该形式空间中的想法。
这是当前实现的一个明确限制：单靠提示词难以支撑更长的运行。

\subsection{消融实验}

要判断模型的贡献，需要与不使用模型的对照组比较。
对照组的做法是：从同一形式空间中均匀随机抽取同样条数的形式，
其余环节（线性规划求解、反例搜索、置换检验、BH 校正）完全相同。
三组随机种子各跑四轮，结果如表\ref{tab:ablation} 所示。

\begin{table}[H]
\centering
\small
\caption{模型提议与随机提议的对比（三组种子的均值）}
\label{tab:ablation}
\begin{tabular}{@{}lccc@{}}
\toprule
指标 & 模型提议 & 随机提议 & 逐种子比较 \\
\midrule
非平凡候选 & 10.67 & 11.33 & 模型 0 胜 1 平 2 负 \\
通过反驳 & 4.33 & 7.00 & 模型 0 胜 3 负 \\
通过筛选 & 2.33 & 1.33 & 模型 2 胜 1 负 \\
测试集成立 & 2.33 & 1.33 & 模型 2 胜 1 负 \\
\bottomrule
\end{tabular}
\end{table}

\begin{figure}[H]
\centering
\includegraphics[width=0.68\linewidth]{fig12_ablation.png}
\caption{模型提议与随机提议在各指标上的对比}
\label{fig:ablation}
\end{figure}

前两项随机组更高，后两项模型组更高，方向相反。
这一结果需要分开解释。

\textbf{随机组在前两项领先，原因在于恒真式。}
九个不变量之间存在由定义直接给出的序关系，
例如 $\omega \leq \chi$、$\delta \leq \Delta$、$\rho \leq \Delta$。
随机抽到的形式有相当比例属于这一类，它们在任意图上都成立，
既不会被反例推翻，也必然通过置换检验。
换言之，随机组在前两关的高分来自“提不出错的东西”，而非“提出了有价值的东西”。

\textbf{模型组在后两项反超，原因在于取等比例。}
把“通过筛选”除以“通过反驳”，得到条件通过率：
模型组 $2.33/4.33 = 54\%$，随机组 $1.33/7.00 = 19\%$。
模型提出的形式一旦躲过反例，就有过半的机会同时通过统计检验；
随机组的形式大多在最后一关被淘汰。
这是因为置换检验的统计量是验证集上的取等比例，
而模型倾向于提出结构上更“紧”的不等式，取等比例天然更高。

\textbf{这一结论并不稳固。}逐种子看，第三组中随机组反而保留了 3 条、
模型组只保留 2 条；三组种子下模型组胜两次、负一次，
不足以支撑“模型优于随机”的判断。
合理的表述是：在本课题的规模下，模型提议在质量上有微弱优势的迹象，
但样本量不足以确认，也无法排除这一差异来自随机波动。

\textbf{更值得注意的是结论随检查器的改进而翻转。}
第一次消融时随机组在全部四项指标上大幅领先，
检查保留清单后才发现其中多是 $\omega \leq \chi$、$n \geq \chi$ 这样的恒真式。
据此增加了恒真式识别模块：用基本事实（定义、简单计数、经典结论）
判断候选是否被其中某一条处处支配，若是则判定为无信息量并剔除。
模型组与随机组此后使用同一剔除规则。
第\ref{sec:ceiling}节又发现该模块漏掉了同一不等式的对称写法，
把基本事实从 32 条补到 35 条，随机组的保留数随之从 4.67 降到 1.33。
两次修正都只影响随机组，因为模型很少提出这类平凡形式。
这说明消融实验的结论对“什么算有信息量”这一判定高度敏感，
而该判定本身需要一个逐步完善的过程。
把任何一次消融的数字当作定论都是不严谨的。

\subsection{穷举上界：智能体离最好结果有多远}
\label{sec:ceiling}

消融实验回答的是“模型是否优于随机”，但没有回答另一个问题：
智能体五轮只提出 40 条，而同一形式空间里究竟有多少条能通过全部检验？
若穷举也只留下四五条，则智能体已接近上限；若穷举能留下几十条，
则它漏掉了很多。为此对完整形式空间逐条求解，施加与智能体相同的四道关卡。

结果如表\ref{tab:ceiling}。形式空间共 2304 条（9 个目标、预测子不超过 2 个、双向），
逐层筛选后只剩 34 条。第三关与智能体不同：智能体只能用有限预算搜索反例，
这里直接判定候选在验证集 $V$ 上是否处处成立，因此 34 是\textbf{上界}，
智能体的产出不可能超过它。

\begin{table}[H]
\centering
\small
\caption{穷举上界的逐层筛选}
\label{tab:ceiling}
\begin{tabularx}{\linewidth}{@{}>{\raggedright\arraybackslash}p{5.4cm}
                              >{\raggedleft\arraybackslash}p{2.2cm}
                              >{\raggedright\arraybackslash}X@{}}
\toprule
关卡 & 剩余条数 & 说明 \\
\midrule
形式空间 & 2304 & 9 个目标 $\times$ 预测子 $\leq 2$ $\times$ 双向 \\
非平凡（线性规划未压成零） & 1236 & 求解耗时 175 秒 \\
按不等式去重 & 358 & 不同形式可拟出同一条 \\
有信息量（不被基本事实支配） & 290 & 剔除 68 条恒真式 \\
在 $V$ 上无反例 & 116 & 智能体此关靠搜索，故此处为上界 \\
通过置换检验（不校正，$p \leq 0.05$） & 35 & 淘汰 81 条 \\
通过 BH 校正（$q = 0.1$） & 34 & 再淘汰 1 条 \\
\bottomrule
\end{tabularx}
\end{table}

三点观察。

第一，智能体的 3 条全部落在上界的 34 条之内，排名分别为第 1、第 7、第 12
（按验证集取等比例排序）。这说明它提出的候选确实属于
“能被穷举认可”的那一类，没有提出无意义的形式。

第二，本节的第一次运行给出过一个看似更有说服力的数字：
上界中取等比例最高的是一条 $n \geq 2\nu$（0.5854），而智能体五轮四十条从未提出。
初看这是提示词引导模型只关注色数、独立数等“难”量所致。
但核对后发现问题不在模型：$n \geq 2\nu$ 与 $n \geq \alpha+\delta$ 之类的形式，
其实是 $\nu \leq n/2$ 的移项，而 $\nu \leq n/2$ 本就在 32 条基本事实之中。
恒真式识别模块只比较目标与方向都相同的写法，
故这类对称改写逃过了检查，被当成有信息量的候选。
补上对称写法之后（基本事实由 32 条增至 35 条），
$n \geq 2\nu$ 被正确剔除，上界从 37 条降到 34 条。
换言之，此前那条“最高分候选”并不构成智能体的漏项，而是检查器的一个漏洞。

第三，也是本节最重要的一点：\textbf{置换检验淘汰了 81 条，BH 校正只淘汰了 1 条}。
这与此前安慰剂实验的结论一致——校正环节在当前规模下几乎不起作用。
原因可以从数据中直接读出：34 条候选里有 33 条的 $p$ 值等于下限 0.004975。
$B = 200$ 时 $p$ 只能取 $1/(B+1)$ 的整数倍，零分布从未达到观测值，
所有强候选的 $p$ 值便挤在同一个下限上，BH 排序时无法区分，自然一条也不删。

这不是校正方法失效，而是统计量在当前分辨率下饱和。
要让校正在这一规模上产生区分，需要 $B$ 至少达到数千，
使 $p$ 值下限降到 $10^{-3}$ 以下。但 $B = 2000$ 时单条形式的置换耗时约 60 秒，
358 条候选需十余小时，超出课程周期的计算预算。
这是本课题一项明确的限制，第\ref{sec:evaluation}章据此讨论了校正的适用条件。

\subsection{本章小结}

本章报告了十组结果，可分为两部分。

前六节检验确定性骨架。数据集的四项已知结论零违反，
线性规划能复现五条已知界，三类反例搜索策略各有其有效区间，
统计筛选能识别出一条过拟合候选。
其中第\ref{sec:injection}节的注入实验给出一条负面结论：
在反例稀少的难例上，三类策略的检出率低至 0.02，
且几乎不随预算改善。这使“未找到反例”这一判断的可信度大幅下降。

后四节接入模型并检验其贡献。系统能产出在测试集上成立且非平凡的候选，
三条保留候选全部落在穷举上界的三十四条之内，其中一条排名第一。
但消融实验给出的证据不足以确认模型优于随机抽样：
模型组在后两项指标上领先，却只在三组种子中的两组取胜，
且这一领先是在两轮检查器修正之后才出现的。

本章另有一项结果属于方法论层面。
两轮消融得到的方向相反的结论，其差异完全来自恒真式识别模块的完善程度。
第一次消融显示随机组全面领先，补上对称写法后又变为模型组在后两项反超。
两次修正都只影响随机组。这说明在自动猜想这类任务中，
“什么算有信息量”这一判定本身是实验的一部分，
它随检查器的改进而移动，因而任何一次消融的数字都只是阶段性结论。
这一教训比消融的具体数值更有价值，第\ref{sec:evaluation}章据此重述了评价口径。
